%========================================================= SECTION 3: PORT SCANNING & BANNER GRABBING
\section{Phase 1 \& 2: TCP SYN Port Scanning and Service Banner Grabbing}\label{sec:scan_banner}

Network reconnaissance begins by identifying active listening transport endpoints and interrogating
exposed services for daemon identity strings. This section details the protocol handshakes, timing
mechanics, and console logs captured during scanning.

\subsection{Phase 1: TCP SYN Half-Open Port Scanning Methodology}
Standard TCP socket connections conform to RFC~793 by completing a full three-way handshake:
the client transmits a SYN packet, the server replies with SYN/ACK, and the client returns an ACK.
Application daemons commit connections to access logs only after this third ACK arrives.

In contrast, our Layer-2 scanner executes a \textbf{half-open stealth scan} operating below the
operating system socket abstraction:
\begin{itemize}
    \item \textbf{Open Port ($\text{SYN}=1, \text{ACK}=1$):} When a service is listening, the target
    kernel returns a SYN/ACK frame ($\text{Flags} = \mathtt{0x12}$). The scanner records the port as
    \textsc{open}, ingests the initial TTL, Window Size, and TCP Options into the fingerprinting
    telemetry buffer, and \hl{deliberately withholds the final ACK}. The TCP state machine on the
    target host eventually expires the half-open embryonic connection via standard timeout, ensuring
    the probe is never committed to application-level access logs.
    \item \textbf{Closed Port ($\text{RST}=1$):} When no daemon listens on the destination port, the
    target kernel emits a TCP RST frame ($\text{Flags} = \mathtt{0x04}$, Window$=0$). The scanner
    marks the port as \textsc{closed} and routes the frame to the \textbf{Closed-Port RST Fallback
    Pipeline}.
    \item \textbf{Filtered Port (Packet Drop / Timeout):} If a network firewall or host filtering
    agent drops the probe, no frame returns. The scanner maintains a non-blocking select loop with a
    calibrated timeout ($t_{\text{timeout}} = 1.5$\,s) before classifying the port as \textsc{filtered}.
\end{itemize}

\begin{figure}[htbp]
\centering
\begin{minipage}[t]{0.238\textwidth}
\centering
{\color{accent}\sffamily\bfseries\scriptsize (a) RFC 793 Handshake}\\[2pt]
\resizebox{\linewidth}{!}{%
\begin{tikzpicture}[font=\sffamily\scriptsize]
    \node[draw=accent, fill=accentsoft, rounded corners=2pt, font=\bfseries, inner sep=2.5pt, minimum width=1.4cm] (c) at (0, 0) {Client};
    \node[draw=rule, fill=codebg, rounded corners=2pt, font=\bfseries, inner sep=2.5pt, minimum width=1.4cm] (s) at (2.6, 0) {Server};
    \draw[thick, color=accent] (c.south) -- (0, -3.3);
    \draw[thick, color=rule] (s.south) -- (2.6, -3.3);
    \draw[-{Stealth[length=4pt]}, thick] (0, -0.5) -- node[above, sloped, font=\tiny] {SYN ($seq = x$)} (2.6, -1.0);
    \draw[-{Stealth[length=4pt]}, thick] (2.6, -1.4) -- node[above, sloped, font=\tiny] {SYN/ACK ($ack = x+1$)} (0, -1.9);
    \draw[-{Stealth[length=4pt]}, thick] (0, -2.3) -- node[above, sloped, font=\tiny] {ACK ($ack = y+1$)} (2.6, -2.8);
    \node[draw=accent, fill=accentsoft, rounded corners=2pt, font=\tiny\bfseries, align=center] at (1.3, -3.1) {ESTABLISHED};
\end{tikzpicture}%
}
\end{minipage}%
\hfill
\begin{minipage}[t]{0.238\textwidth}
\centering
{\color{accent}\sffamily\bfseries\scriptsize (b) SYN Scan: Open}\\[2pt]
\resizebox{\linewidth}{!}{%
\begin{tikzpicture}[font=\sffamily\scriptsize]
    \node[draw=accent, fill=accentsoft, rounded corners=2pt, font=\bfseries, inner sep=2.5pt, minimum width=1.4cm] (sc) at (0, 0) {Scanner};
    \node[draw=rule, fill=codebg, rounded corners=2pt, font=\bfseries, inner sep=2.5pt, minimum width=1.4cm] (th) at (2.6, 0) {Target};
    \draw[thick, color=accent] (sc.south) -- (0, -3.3);
    \draw[thick, color=rule] (th.south) -- (2.6, -3.3);
    \draw[-{Stealth[length=4pt]}, thick] (0, -0.5) -- node[above, sloped, font=\tiny] {Raw SYN (:80)} (2.6, -1.0);
    \draw[-{Stealth[length=4pt]}, thick, color=accent] (2.6, -1.3) -- node[above, sloped, font=\tiny] {SYN/ACK (\texttt{0x12})} (0, -1.8);
    \draw[-{Stealth[length=4pt]}, thick, dashed, color=accentstrong] (0, -2.3) -- (1.6, -2.3);
    \draw[thick, color=accentstrong] (1.5, -2.15) -- (1.7, -2.45);
    \draw[thick, color=accentstrong] (1.5, -2.45) -- (1.7, -2.15);
    \node[above=1pt, font=\tiny\color{accentstrong}\bfseries] at (0.9, -2.25) {Withhold ACK};
    \node[draw=accent, fill=accentsoft, rounded corners=2pt, font=\tiny\bfseries, align=center] at (1.3, -3.1) {State: \textbf{OPEN}};
\end{tikzpicture}%
}
\end{minipage}%
\hfill
\begin{minipage}[t]{0.238\textwidth}
\centering
{\color{accent}\sffamily\bfseries\scriptsize (c) SYN Scan: Closed}\\[2pt]
\resizebox{\linewidth}{!}{%
\begin{tikzpicture}[font=\sffamily\scriptsize]
    \node[draw=accent, fill=accentsoft, rounded corners=2pt, font=\bfseries, inner sep=2.5pt, minimum width=1.4cm] (sc) at (0, 0) {Scanner};
    \node[draw=rule, fill=codebg, rounded corners=2pt, font=\bfseries, inner sep=2.5pt, minimum width=1.4cm] (th) at (2.6, 0) {Target};
    \draw[thick, color=accent] (sc.south) -- (0, -3.3);
    \draw[thick, color=rule] (th.south) -- (2.6, -3.3);
    \draw[-{Stealth[length=4pt]}, thick] (0, -0.5) -- node[above, sloped, font=\tiny] {Raw SYN (:443)} (2.6, -1.0);
    \draw[-{Stealth[length=4pt]}, thick, color=accentstrong] (2.6, -1.4) -- node[above, sloped, font=\tiny] {RST (\texttt{0x04})} (0, -1.9);
    \node[draw=rule, fill=codebg, rounded corners=2pt, font=\tiny\bfseries, align=center] at (1.3, -2.8) {State: \textbf{CLOSED}};
\end{tikzpicture}%
}
\end{minipage}%
\hfill
\begin{minipage}[t]{0.238\textwidth}
\centering
{\color{accent}\sffamily\bfseries\scriptsize (d) SYN Scan: Filtered}\\[2pt]
\resizebox{\linewidth}{!}{%
\begin{tikzpicture}[font=\sffamily\scriptsize]
    \node[draw=accent, fill=accentsoft, rounded corners=2pt, font=\bfseries, inner sep=2.5pt, minimum width=1.4cm] (sc) at (0, 0) {Scanner};
    \node[draw=rule, fill=codebg, rounded corners=2pt, font=\bfseries, inner sep=2.5pt, minimum width=1.4cm] (th) at (2.6, 0) {Firewall};
    \draw[thick, color=accent] (sc.south) -- (0, -3.3);
    \draw[thick, color=rule] (th.south) -- (2.6, -3.3);
    \draw[-{Stealth[length=4pt]}, thick] (0, -0.5) -- node[above, sloped, font=\tiny] {Raw SYN (:8080)} (2.3, -0.9);
    \node[draw=accentstrong, fill=accentsoft, rounded corners=1pt, font=\tiny\bfseries, text=accentstrong] at (2.6, -0.9) {DROP};
    \draw[thick, dashed, color=rule] (0.4, -1.6) -- (2.2, -1.6) node[midway, above, font=\tiny] {No Response};
    \node[draw=rule, fill=codebg, rounded corners=2pt, font=\tiny, align=center] at (1.3, -2.1) {Wait $1.5\,\text{s}$};
    \node[draw=rule, fill=codebg, rounded corners=2pt, font=\tiny\bfseries, align=center] at (1.3, -2.8) {State: \textbf{FILTERED}};
\end{tikzpicture}%
}
\end{minipage}
\caption{Protocol Timing Diagrams: RFC~793 3-Way Handshake vs.\ Half-Open Probe Classifications.}
\label{fig:timing_scan}
\end{figure}

\FloatBarrier
\subsubsection{Empirical Evaluation: 100-Port Comprehensive Range Scan on Windows Target}
To evaluate scalability, port diversity, and safety constraints ($\text{MAX\_PROBES} = 50$), the
engine executed a 100-port scan against a live Windows 11 target (\texttt{192.168.1.104}) in two
consecutive 50-port batches:

\begin{logbox}
=============================================================================
100-PORT COMPREHENSIVE SCALE VERIFICATION RUN
Target IP: 192.168.1.104 | Attacker IP: 192.168.1.105 (d8:80:83:74:7c:b3)
Target MAC: 98:fe:3e:5b:c3:27 | Inter-Probe Delay: 0.50s
=============================================================================

=== BATCH 1: Ports 20 to 1723 (Top 50 Well-Known Networking & System Services) ===
[PHASE 1] TCP SYN Port Scanning:
    20  CLOSED                         TTL=128 WIN=0
    21  CLOSED                         TTL=128 WIN=0
    22  CLOSED                         TTL=128 WIN=0
    23  CLOSED                         TTL=128 WIN=0
    25  CLOSED                         TTL=128 WIN=0
    53  CLOSED                         TTL=128 WIN=0
    80  CLOSED                         TTL=128 WIN=0
   135  OPEN                           TTL=128 WIN=65535
   137  CLOSED                         TTL=128 WIN=0
   138  CLOSED                         TTL=128 WIN=0
   139  OPEN                           TTL=128 WIN=8192
   161  CLOSED                         TTL=128 WIN=0
   443  CLOSED                         TTL=128 WIN=0
   445  OPEN                           TTL=128 WIN=65535
  1025  OPEN                           TTL=128 WIN=65535
  1433  CLOSED                         TTL=128 WIN=0
  1723  CLOSED                         TTL=128 WIN=0
[Summary: 4 Open, 46 Closed, 0 Filtered]

[PHASE 3] OS Fingerprinting (Candidate: Open-Port TCP SYN/ACK on Port 135)
  Likely OS  : Windows                 Confidence : 71%
  Version    : Windows 10 / 11 / Server 2016+
               71% Windows 10 / 11 / Server 2016+
               29% Windows 7 / 8 / Vista / Server 2008-2012
  TTL        : 128 (initial 128, 0 hops)   Window: 65535    DF bit: set
  Options    : ('MSS', 1460), ('NOP',), ('WSCALE', 8), ('SACK_PERMITTED',), ('TIMESTAMP', ...)
  Evidence   : initial TTL 128 | Window 65535 | Windows Option Order | IP-ID strictly incremental (deltas: [1, 1, 1, ...], avg 1.3) | Timestamp clock ~1001.0 Hz (1ms tick)

=== BATCH 2: Ports 1883 to 49164 (Top 50 High Ports, RPC, Cloud & Web Services) ===
[PHASE 1] TCP SYN Port Scanning (Evaluating Zero-Open-Port Environment):
  1883  CLOSED  TTL=128 WIN=0    2049  CLOSED  TTL=128 WIN=0
  2121  CLOSED  TTL=128 WIN=0    2221  CLOSED  TTL=128 WIN=0
  3000  CLOSED  TTL=128 WIN=0    3306  CLOSED  TTL=128 WIN=0
  3389  CLOSED  TTL=128 WIN=0    5000  CLOSED  TTL=128 WIN=0
  8080  CLOSED  TTL=128 WIN=0    8443  CLOSED  TTL=128 WIN=0
  9000  CLOSED  TTL=128 WIN=0   49152  CLOSED  TTL=128 WIN=0
[Summary: 0 Open, 50 Closed, 0 Filtered]

[PHASE 3] OS Fingerprinting (Candidate: Closed-Port TCP RST Fallback)
  Likely OS  : Windows                 Confidence : 69%
  Version    : Windows family (closed-port RST fallback telemetry)
  TTL        : 128 (initial 128, 0 hops)   Window: 0        DF bit: set
  IP ID Seq  : [13207, 13208, 13209, 13210, 13211, 13212, 13213, 13215, 13217, 13218...]
  IP ID Class: STRICT_INCREMENTAL (deltas: [1, 1, 1, 1, 1, 1, 2, 2, 1, 1...], avg 1.1)
  Evidence   : TCP RST initial TTL 128 | Window 0 | Strict incremental IP ID global counter | Zero open ports required!
\end{logbox}

\subsection{Phase 2: Service Banner Grabbing Methodology}
For every open port discovered during Phase~1, the engine initiates a stream socket handshake to
capture software identity strings:
\begin{itemize}
    \item \textbf{Passive Services (FTP, SMTP):} Listening daemons transmit an unprompted banner
    immediately following the three-way handshake. The grabber ingests the banner buffer and exits cleanly.
    \item \textbf{Active Services (HTTP):} Web daemons remain silent until an HTTP verb is received.
    The grabber establishes a standard connection, transmits an RFC-compliant probe (\texttt{HEAD / HTTP/1.0}),
    and parses the \texttt{Server:} header line.
\end{itemize}

\begin{figure}[htbp]
\centering
\begin{minipage}[t]{0.485\textwidth}
\centering
{\color{accent}\sffamily\bfseries\scriptsize (a) Passive Service Banner: FTP (Port 21)}\\[4pt]
\begin{tikzpicture}[font=\sffamily\scriptsize, scale=0.78, every node/.style={transform shape}]
    \node[draw=accent, fill=accentsoft, rounded corners=2pt, font=\bfseries, inner sep=3pt, minimum width=1.8cm] (g) at (0, 0) {Banner Grabber};
    \node[draw=rule, fill=codebg, rounded corners=2pt, font=\bfseries, inner sep=3pt, minimum width=1.8cm] (s) at (4.0, 0) {FTP Server (:21)};
    \draw[thick, color=accent] (g.south) -- (0, -3.8);
    \draw[thick, color=rule] (s.south) -- (4.0, -3.8);
    \draw[-{Stealth[length=4pt]}, thick] (0, -0.5) -- node[above, sloped, font=\tiny] {TCP 3-Way Handshake} (4.0, -1.0);
    \draw[-{Stealth[length=4pt]}, thick, color=accent] (4.0, -1.4) -- node[above, sloped, font=\tiny\ttfamily] {220 LabFTP Ready\textbackslash r\textbackslash n} (0, -1.9);
    \draw[-{Stealth[length=4pt]}, thick] (0, -2.3) -- node[above, sloped, font=\tiny\ttfamily] {QUIT\textbackslash r\textbackslash n} (4.0, -2.8);
    \draw[-{Stealth[length=4pt]}, thick] (4.0, -3.0) -- node[above, sloped, font=\tiny\ttfamily] {221 Goodbye} (0, -3.4);
\end{tikzpicture}
\end{minipage}%
\hfill
\begin{minipage}[t]{0.485\textwidth}
\centering
{\color{accent}\sffamily\bfseries\scriptsize (b) Active Service Banner: HTTP (Port 80)}\\[4pt]
\begin{tikzpicture}[font=\sffamily\scriptsize, scale=0.78, every node/.style={transform shape}]
    \node[draw=accent, fill=accentsoft, rounded corners=2pt, font=\bfseries, inner sep=3pt, minimum width=1.8cm] (g) at (0, 0) {Banner Grabber};
    \node[draw=rule, fill=codebg, rounded corners=2pt, font=\bfseries, inner sep=3pt, minimum width=1.8cm] (s) at (4.0, 0) {HTTP Server (:80)};
    \draw[thick, color=accent] (g.south) -- (0, -3.8);
    \draw[thick, color=rule] (s.south) -- (4.0, -3.8);
    \draw[-{Stealth[length=4pt]}, thick] (0, -0.5) -- node[above, sloped, font=\tiny] {TCP 3-Way Handshake} (4.0, -1.0);
    \draw[-{Stealth[length=4pt]}, thick] (0, -1.4) -- node[above, sloped, font=\tiny\ttfamily] {HEAD / HTTP/1.0\textbackslash r\textbackslash n\textbackslash r\textbackslash n} (4.0, -1.9);
    \draw[-{Stealth[length=4pt]}, thick, color=accent] (4.0, -2.3) -- node[above, sloped, font=\tiny\ttfamily] {Server: VulnHTTP/1.0} (0, -2.8);
    \draw[-{Stealth[length=4pt]}, thick] (0, -3.1) -- node[above, sloped, font=\tiny] {TCP FIN Close} (4.0, -3.5);
\end{tikzpicture}
\end{minipage}
\caption{Timing Sequences for Passive Banner Capture and Active Application Probing.}
\label{fig:timing_banner}
\end{figure}

\FloatBarrier
\subsection{Empirical Service Banner Telemetry Across Evaluation Targets}
The Phase~2 banner output from each testbed target is shown below:

\begin{logbox}
=================================================================
  [PHASE 2] Service Banner Grabbing: Windows Lab Host (192.168.0.116)
=================================================================
  Target      : 192.168.0.116 (Windows 11 Lab Server)
  Probed Ports: 21 (FTP), 80 (HTTP)
  Timeout     : 2.0s per socket
-----------------------------------------------------------------
  [OK]    21  FTP      -> LabFTP Server (Vulnerable Demo) Ready
  [OK]    80  HTTP     -> VulnHTTP/1.0 (Lab Demo)
-----------------------------------------------------------------
  Banners found : 2/2 (100% Acquisition Rate)
  Cross-Layer   : Correlated daemon signatures dispatched to Phase 3
=================================================================

=================================================================
  [PHASE 2] Service Banner Grabbing: Kali Linux Target (192.168.56.101)
=================================================================
  Target      : 192.168.56.101 (Debian GNU/Linux Kernel 6.x VM)
  Probed Ports: 22 (SSH), 80 (HTTP), 8080 (HTTP)
  Timeout     : 2.0s per socket
-----------------------------------------------------------------
  [OK]    22  SSH      -> SSH-2.0-OpenSSH_10.3p1 Debian-5
  [OK]    80  HTTP     -> Apache/2.4.68 (Debian)
  [OK]  8080  HTTP     -> SimpleHTTP/0.6 Python/3.13.14
-----------------------------------------------------------------
  Cross-Layer: 4 banner signature(s) correlated
    - Cross-layer banner on port 22 (SSH) matches Linux signature 'Debian' (Debian GNU/Linux distribution)
    - Cross-layer banner on port 80 (HTTP) matches Linux signature 'Debian' (Debian GNU/Linux distribution)
  Banners found : 3/3 (100% Acquisition Rate)
=================================================================
\end{logbox}

\begin{logbox}
=================================================================
  [PHASE 2] Service Banner Grabbing: Windows OpenSSH Target (192.168.1.104)
=================================================================
  Target      : 192.168.1.104 (AisheePC Windows 11 Desktop)
  Probed Ports: 22 (SSH)
  Timeout     : 2.0s per socket
-----------------------------------------------------------------
  [OK]    22  SSH      -> SSH-2.0-OpenSSH_for_Windows_8.1
-----------------------------------------------------------------
  Cross-Layer: 1 banner signature(s) correlated
    - Cross-layer banner on port 22 (SSH) matches Windows signature 'OpenSSH_for_Windows' (Microsoft Windows OpenSSH port)
  Banners found : 1/1 (100% Acquisition Rate)
=================================================================

=================================================================
  [PHASE 2] Service Banner Grabbing: Mobile Subnet Targets (iPad & Android)
=================================================================
  [Target Node: 192.168.0.152 (Apple iPad)]
  Probed Ports: 80 (HTTP)
  [OK]    80  HTTP     -> HTTP/1.1 200 OK | Date: Fri, 11 Sep 2026 21:15:02 GMT
                          Server: VLC/3.0.19 LibVLC/3.0.19 (Media Sharing Web Server)
  Banners found : 1/1 (100% Acquisition Rate)

  [Target Node: 192.168.0.169 (Android 14 Phone)]
  Probed Ports: 2221 (FTP)
  [OK]  2221  FTP      -> 220 Service ready for new user. (WiFi FTP Server Daemon)
  Banners found : 1/1 (100% Acquisition Rate)
=================================================================
\end{logbox}
